\documentclass[12pt, a4paper]{article}
% --- 1. Cấu hình Tiếng Việt và Font chữ chuẩn ---
\usepackage[utf8]{inputenc}
\usepackage[T5]{fontenc}
\usepackage[vietnamese]{babel}
\usepackage{mathptmx} % Font Times New Roman đồng bộ cả chữ và công thức toán, không gây xung đột gói

\makeatletter
\renewcommand{\normalsize}{\@setfontsize\normalsize{13pt}{16pt}}
\makeatother
\normalsize

% --- 2. Gói Toán học & Ký hiệu bổ trợ ---
\usepackage{amsmath, amssymb, amsfonts, mathrsfs, amsthm}
\usepackage{graphicx}
\usepackage{fontawesome}
\usepackage{booktabs, tabularx, array, multirow}
\usepackage{enumitem}

% --- 3. Đồ họa TikZ, Bảng biến thiên & Hình không gian ---
\usepackage{tikz}
\usetikzlibrary{calc, intersections, angles, quotes, patterns, arrows.meta, 3d}
\usepackage{pgfplots}
\pgfplotsset{compat=1.18}
\usepackage{tkz-euclide}
\usepackage{tkz-tab}

\renewcommand{\vec}[1]{\overrightarrow{#1}}

% --- 4. Hộp sư phạm (Tcolorbox) ---
\usepackage[many]{tcolorbox}
\tcbuselibrary{skins, breakable}

% --- 5. Căn lề chuẩn văn bản giáo án ---
\usepackage{geometry}
\geometry{
    a4paper,
    left=25mm,
    right=15mm,
    top=20mm,
    bottom=20mm
}

% --- 6. Header / Footer chuẩn hóa ---
\usepackage{fancyhdr}
\usepackage{lastpage}
\pagestyle{fancy}
\fancyhf{}

\fancyhead[L]{\small \textit{Kế hoạch bài dạy Toán 11}}
\fancyhead[R]{\textbf{Trang \thepage/\pageref{LastPage}}}
\fancyfoot[L]{\small \textit{Tổ Toán - Tin}}
\fancyfoot[R]{\small \textit{Trường THCS\&THPT Hồng Vân}}
\renewcommand{\headrulewidth}{0.4pt}
\renewcommand{\footrulewidth}{0.4pt}

% --- 7. Giãn dòng & Giãn đoạn theo chuẩn Nghị định 30/2020/NĐ-CP ---
\usepackage{setspace}
\setstretch{1.15}
\setlength{\parindent}{1.0cm}
\setlength{\parskip}{5pt}

% Định nghĩa hộp sư phạm chuyên dụng
\newtcolorbox{pedabox}[2][]{
    enhanced,
    breakable,
    colback=blue!3!white,
    colframe=blue!65!black,
    fonttitle=\bfseries\small,
    title={#2},
    arc=2mm,
    boxrule=0.6pt,
    left=3mm, right=3mm, top=2mm, bottom=2mm,
    #1
}

\newtcolorbox{aibox}[2][]{
    enhanced,
    breakable,
    colback=teal!4!white,
    colframe=teal!70!black,
    fonttitle=\bfseries\small,
    title={#2},
    arc=2mm,
    boxrule=0.6pt,
    left=3mm, right=3mm, top=2mm, bottom=2mm,
    #1
}

\newtcolorbox{scaffoldbox}[2][]{
    enhanced,
    breakable,
    colback=orange!4!white,
    colframe=orange!80!black,
    fonttitle=\bfseries\small,
    title={#2},
    arc=2mm,
    boxrule=0.6pt,
    left=3mm, right=3mm, top=2mm, bottom=2mm,
    #1
}

\newtcolorbox{stembox}[2][]{
    enhanced,
    breakable,
    colback=purple!4!white,
    colframe=purple!75!black,
    fonttitle=\bfseries\small,
    title={#2},
    arc=2mm,
    boxrule=0.6pt,
    left=3mm, right=3mm, top=2mm, bottom=2mm,
    #1
}

\begin{document}

\begin{center}
    {\fontsize{14pt}{17pt}\selectfont \textbf{KẾ HOẠCH BÀI DẠY MÔN TOÁN LỚP 11}}\\[6pt]
    {\fontsize{16pt}{19pt}\selectfont \textbf{BÀI 21. PHƯƠNG TRÌNH, BẤT PHƯƠNG TRÌNH MŨ VÀ LÔGARIT}}\\[4pt]
    {\fontsize{12pt}{15pt}\selectfont \textbf{Môn học: Toán 11} \ \textbar \ \textbf{Thời lượng: 2 tiết (Tiết 60, 61 theo PPCT | Tuần 20)}}
\end{center}

\vspace{5pt}

\section*{I. MỤC TIÊU}

\subsection*{1. Kiến thức, Năng lực}

\begin{itemize}[leftmargin=1.2cm]
    \item \textbf{Yêu cầu cần đạt (Nguyên văn CT GDPT 2018 Toán 11):}
    \begin{itemize}[leftmargin=0.6cm]
        \item Giải được phương trình, bất phương trình mũ, lôgarit ở dạng đơn giản (ví dụ: $2^{x+1} = \dfrac{1}{4}$; $2^{x+1} = 2^{3x+5}$; $\log_2(x+1) = 3$; $\log_3(x+1) = \log_3(x^2-1)$).
        \item Giải quyết được một số vấn đề có liên quan đến môn học khác hoặc có liên quan đến thực tiễn gắn với phương trình, bất phương trình mũ và lôgarit (ví dụ: bài toán xác định niên đại khảo cổ bằng đồng vị phóng xạ Carbon-14, bài toán tính độ rung chấn động đất theo thang đo Richter, bài toán kiểm soát độ pH trong xử lý nước).
    \end{itemize}

    \item \textbf{Năng lực Toán học đặc thù:}
    \begin{itemize}[leftmargin=0.6cm]
        \item \textit{Năng lực tư duy và lập luận toán học:} So sánh, phân tích bản chất cơ số $a > 1$ (bất phương trình giữ nguyên chiều) và $0 < a < 1$ (bất phương trình bắt buộc đổi chiều); lập luận chặt chẽ về điều kiện xác định của biểu thức trong dấu lôgarit trước khi biến đổi tương đương; nhận diện các nghiệm ngoại lai (spurious roots).
        \item \textit{Năng lực mô hình hóa toán học:} Chuyển đổi các bài toán thực tiễn (phân rã phóng xạ, tăng trưởng tài chính, độ Richter địa chấn) thành mô hình phương trình, bất phương trình mũ và lôgarit để tìm ẩn số thời gian hoặc năng lượng.
        \item \textit{Năng lực giải quyết vấn đề toán học:} Lựa chọn và thực hiện phương pháp giải phù hợp: đưa về cùng cơ số, lôgarit hóa / mũ hóa hai vế, đặt ẩn phụ quy về phương trình bậc hai; đối chiếu tập nghiệm với điều kiện xác định ban đầu.
        \item \textit{Năng lực giao tiếp toán học:} Trình bày lời giải phương trình, bất phương trình rõ ràng, mạch lạc, đúng chuẩn mực ký hiệu logic toán học ($\iff, \implies, \in, \cap$); diễn đạt chính xác tập nghiệm bằng khoảng, đoạn hoặc hợp các khoảng.
        \item \textit{Năng lực sử dụng công cụ, phương tiện học toán:} Sử dụng thành thạo máy tính cầm tay (MTCT) Casio fx-580VN X: chức năng \texttt{SHIFT SOLVE} để dò tìm nghiệm phương trình mũ và lôgarit; sử dụng tính năng \texttt{CALC} để thử lại nghiệm và kiểm tra dấu bất phương trình trên từng khoảng phân định.
    \end{itemize}

    \item \textbf{Năng lực số (Theo Thông tư 02/2025/TT-BGDĐT):}
    \begin{itemize}[leftmargin=0.6cm]
        \item \textit{Mã 5.2NC1b (Giải quyết vấn đề bằng công cụ toán học số):} Học sinh sử dụng máy tính cầm tay thực hiện lệnh \texttt{SOLVE} để tìm nghiệm xấp xỉ của các phương trình mũ, lôgarit; biết phân tích sai số giải thuật lặp và ý nghĩa giá trị ban đầu (Start value) khi máy tính dò nghiệm phương trình có nhiều nghiệm.
    \end{itemize}

    \item \textbf{Năng lực AI (Theo Quyết định 2422/QĐ-BGDĐT):}
    \begin{itemize}[leftmargin=0.6cm]
        \item \textit{Mã 11.A1.1 (Kiểm chứng và Đánh giá tính chính xác của Trí tuệ nhân tạo):} Học sinh đối chiếu, kiểm chứng các bước biến đổi phương trình do trợ lý AI (Gemini, ChatGPT, PhotoMath) gợi ý; phát hiện các lỗi sai kinh điển của AI khi bỏ quên điều kiện xác định ($f(x) > 0$), dẫn đến việc nhận nghiệm ngoại lai làm sai lệch tập nghiệm của bài toán.
    \end{itemize}

    \item \textbf{Năng lực chung:}
    \begin{itemize}[leftmargin=0.6cm]
        \item \textit{Tự chủ và tự học:} Tự giác hoàn thành phiếu học tập, chủ động phát hiện lỗi sai về đổi chiều bất phương trình và tự sửa chữa.
        \item \textit{Giao tiếp và hợp tác:} Tích cực phối hợp theo cặp đôi để giải quyết các bài toán phân hóa, tranh luận về tính hợp lý của tập nghiệm.
    \end{itemize}
\end{itemize}

\subsection*{2. Phẩm chất}

\begin{itemize}[leftmargin=1.2cm]
    \item \textbf{Chăm chỉ:} Kiên trì thực hiện đầy đủ các bước giải phương trình, không bỏ qua bước đặt điều kiện xác định.
    \item \textbf{Trung thực:} Tự lực biến đổi bài toán, trung thực trong quá trình tự đánh giá và kiểm chứng lời giải với máy tính và AI.
    \item \textbf{Trách nhiệm:} Cẩn thận, chỉn chu trong từng phép biến đổi dấu và phép so sánh số thực; có ý thức ứng dụng toán học vào các vấn đề an toàn thực tiễn (động đất, môi trường).
\end{itemize}

\vspace{5pt}

\section*{II. THIẾT BỊ DẠY HỌC VÀ HỌC LIỆU}

\begin{itemize}[leftmargin=1.2cm]
    \item \textbf{Giáo viên:}
    \begin{itemize}[leftmargin=0.6cm]
        \item Kế hoạch bài dạy chi tiết 2 tiết theo chuẩn Công văn 5512/BGDĐT-GDTrH.
        \item Bài trình chiếu PowerPoint tương tác, tích hợp phần mềm giả lập MTCT Casio fx-580VN X trên máy tính.
        \item Hệ thống Phiếu học tập phân hóa bậc thang: Phiếu học tập số 1 (Tiết 1) và Phiếu học tập số 2 (Tiết 2).
        \item Bảng Rubric đánh giá năng lực học sinh theo chuẩn Thông tư 02/2025 và QĐ 2422.
    \end{itemize}
    \item \textbf{Học sinh:}
    \begin{itemize}[leftmargin=0.6cm]
        \item Sách giáo khoa Toán 11, vở ghi chép bài học, giấy nháp.
        \item Máy tính cầm tay Casio fx-580VN X (hoặc tương đương).
        \item Ôn lại định nghĩa và tính chất của hàm số mũ, hàm số lôgarit đã học ở Bài 20.
    \end{itemize}
\end{itemize}

\vspace{5pt}

\section*{III. TIẾN TRÌNH DẠY HỌC (2 TIẾT | TIẾT 60, 61 THEO PPCT)}

\subsection*{\faClockO \ TIẾT 1 (TIẾT 60 THEO PPCT): PHƯƠNG TRÌNH MŨ VÀ PHƯƠNG TRÌNH LÔGARIT CƠ BẢN}

\subsubsection*{1. Hoạt động 1: Mở đầu (Khởi động) (5 phút)}

\noindent \textbf{a) Mục tiêu:}
Dẫn dắt học sinh nhận biết sự xuất hiện tự nhiên của phương trình mũ từ bài toán phân rã phóng xạ trong vật lí và khảo cổ học, tạo nhu cầu tìm nghiệm của phương trình.

\noindent \textbf{b) Nội dung:}
Giáo viên trình chiếu tình huống:
\begin{pedabox}{Tình huống khởi động: Xác định niên đại cổ vật bằng Carbon-14}
Đồng vị phóng xạ Carbon-14 ($^{14}\text{C}$) có chu kỳ bán rã khoảng $5730$ năm. Khối lượng $^{14}\text{C}$ còn lại trong một mẫu gỗ cổ sau $t$ năm được tính theo công thức:
$$m(t) = m_0 \cdot \left(\dfrac{1}{2}\right)^{\frac{t}{5730}}$$
(trong đó $m_0$ là khối lượng $^{14}\text{C}$ ban đầu lúc sinh vật còn sống).
Một nhà khảo cổ học khai quật được một mẩu than cổ và phát hiện khối lượng $^{14}\text{C}$ chỉ còn lại bằng $\dfrac{1}{8}$ khối lượng ban đầu ($m(t) = \dfrac{1}{8} m_0$).
\begin{enumerate}[leftmargin=0.6cm]
    \item Thiết lập phương trình để tìm số năm $t$ của mẩu than cổ đó.
    \item Hãy tìm giá trị của $t$.
\end{enumerate}
\end{pedabox}

\noindent \textbf{c) Sản phẩm:}
Câu trả lời của học sinh:
\begin{itemize}
    \item Phương trình thiết lập: $m_0 \cdot \left(\dfrac{1}{2}\right)^{\frac{t}{5730}} = \dfrac{1}{8} m_0 \iff \left(\dfrac{1}{2}\right)^{\frac{t}{5730}} = \dfrac{1}{8}$.
    \item Biến đổi: Vì $\dfrac{1}{8} = \left(\dfrac{1}{2}\right)^3$, nên phương trình trở thành:
    $$\left(\dfrac{1}{2}\right)^{\frac{t}{5730}} = \left(\dfrac{1}{2}\right)^3 \iff \dfrac{t}{5730} = 3 \iff t = 3 \times 5730 = 17\,190\text{ năm}.$$
    Mẩu than cổ có niên đại khoảng $17\,190$ năm trước!
\end{itemize}

\noindent \textbf{d) Tổ chức thực hiện:}
\begin{itemize}[leftmargin=0.6cm]
    \item \textbf{Giao nhiệm vụ:} Giáo viên trình chiếu tình huống, yêu cầu học sinh làm việc cá nhân trong 2 phút rồi thảo luận cặp đôi.
    \item \textbf{Thực hiện nhiệm vụ:} Học sinh nhận diện phương trình có ẩn số $t$ nằm ở số mũ, đưa về cùng cơ số $\dfrac{1}{2}$.
    \item \textbf{Báo cáo, thảo luận:} Đại diện 1 học sinh trả lời miệng; cả lớp theo dõi.
    \item \textbf{Kết luận, nhận định:} Giáo viên dẫn dắt vào bài học: Phương trình trên là một ví dụ điển hình của \textbf{phương trình mũ}. Tiếp theo, chúng ta sẽ xây dựng phương pháp tổng quát để giải các phương trình mũ và phương trình lôgarit.
\end{itemize}

\subsubsection*{2. Hoạt động 2: Hình thành kiến thức (25 phút)}

\paragraph{2.1. Đơn vị kiến thức 1: Phương trình mũ cơ bản và phương pháp đưa về cùng cơ số (12 phút)}

\noindent \textbf{a) Mục tiêu:}
Nắm vững định nghĩa và nghiệm của phương trình mũ cơ bản; giải thành thạo phương trình mũ bằng phương pháp đưa về cùng cơ số.

\noindent \textbf{b) Nội dung:}
\begin{pedabox}{Phương trình mũ cơ bản $a^x = b$ ($a > 0, a \neq 1$)}
\begin{itemize}[leftmargin=0.6cm]
    \item Nếu $b \le 0$: Phương trình \textbf{vô nghiệm} (vì $a^x > 0$ với mọi $x \in \mathbb{R}$).
    \item Nếu $b > 0$: Phương trình có \textbf{nghiệm duy nhất} $x = \log_a b$.
\end{itemize}
\end{pedabox}

\begin{pedabox}{Phương pháp đưa về cùng cơ số}
Với $a > 0, a \neq 1$:
$$a^{f(x)} = a^{g(x)} \iff f(x) = g(x)$$
\end{pedabox}

\begin{scaffoldbox}{Giàn giáo sư phạm: Khắc phục sai lầm về vế phải phương trình mũ}
\textbf{Khắc sâu cho học sinh trung bình - yếu:}
\begin{itemize}
    \item Gặp phương trình: $2^x = -4 \implies$ Kết luận ngay phương trình \textbf{VÔ NGHIỆM} (học sinh yếu rất hay nhầm viết $x = -2$).
    \item Gặp phương trình: $3^x = 0 \implies$ Kết luận ngay phương trình \textbf{VÔ NGHIỆM}.
    \item Gặp phương trình: $5^x = 7 \implies$ Nghiệm là $x = \log_5 7$ (không cần tính ra số thập phân trừ khi đề bài yêu cầu).
\end{itemize}
\end{scaffoldbox}

\begin{aibox}{Tích hợp Năng lực số (Mã 5.2NC1b): Bấm phím SOLVE trên MTCT Casio fx-580VN X}
\begin{itemize}[leftmargin=0.6cm]
    \item \textbf{Ví dụ:} Giải phương trình $2^{x+1} = \dfrac{1}{4}$.
    \item \textbf{Thao tác bấm máy:} Nhập biểu thức vào màn hình: \fbox{2} $\to$ \fbox{$x^\blacksquare$} $\to$ \fbox{$\text{ALPHA}$} $\to$ \fbox{)} (nhập biến $X$) $\to$ \fbox{+} $\to$ \fbox{1} $\to$ \fbox{$\blacktriangleright$} $\to$ \fbox{$\text{ALPHA}$} $\to$ \fbox{$\text{CALC}$} (dấu $=$) $\to$ \fbox{$\dfrac{\blacksquare}{\blacksquare}$} $\to$ \fbox{1} $\to$ \fbox{$\blacktriangledown$} $\to$ \fbox{4}.
    \item Bấm lệnh giải: \fbox{$\text{SHIFT}$} $\to$ \fbox{$\text{CALC}$} (\texttt{SOLVE}) $\to$ Máy hỏi \texttt{X=?}, bấm \fbox{0} $\to$ \fbox{$=$}. Màn hình hiển thị: $X = -3$ và $L - R = 0$ (nghiệm chính xác).
\end{itemize}
\end{aibox}

\noindent \textbf{c) Sản phẩm:}
Học sinh thực hiện biến đổi các ví dụ:
\begin{itemize}
    \item \textbf{Ví dụ 1:} $2^{x+1} = \dfrac{1}{4} \iff 2^{x+1} = 2^{-2} \iff x + 1 = -2 \iff x = -3$.
    \item \textbf{Ví dụ 2:} $3^{2x - 1} = 27 \iff 3^{2x - 1} = 3^3 \iff 2x - 1 = 3 \iff 2x = 4 \iff x = 2$.
    \item \textbf{Ví dụ 3:} $2^{x+1} = 2^{3x+5} \iff x + 1 = 3x + 5 \iff 2x = -4 \iff x = -2$.
\end{itemize}

\noindent \textbf{d) Tổ chức thực hiện:}
\begin{itemize}[leftmargin=0.6cm]
    \item \textbf{Giao nhiệm vụ:} Giáo viên trình bày lý thuyết, hướng dẫn giải Ví dụ 1 và thao tác trên MTCT. Yêu cầu học sinh làm Ví dụ 2 và 3 vào vở.
    \item \textbf{Thực hiện nhiệm vụ:} Học sinh biến đổi đưa về cùng cơ số $2$ và $3$; dùng MTCT bấm \texttt{SOLVE} kiểm tra.
    \item \textbf{Báo cáo, thảo luận:} Gọi 2 học sinh lên bảng giải; cả lớp nhận xét.
    \item \textbf{Kết luận, nhận định:} Giáo viên chuẩn hóa các bước giải: Nhận diện cơ số chung $\to$ Biến đổi số mũ $\to$ Giải phương trình đại số.
\end{itemize}

\paragraph{2.2. Đơn vị kiến thức 2: Phương trình lôgarit cơ bản và phương pháp đưa về cùng cơ số (13 phút)}

\noindent \textbf{a) Mục tiêu:}
Nắm vững nghiệm của phương trình lôgarit cơ bản; giải thành thạo phương trình lôgarit bằng phương pháp đưa về cùng cơ số kèm theo điều kiện xác định.

\noindent \textbf{b) Nội dung:}
\begin{pedabox}{Phương trình lôgarit cơ bản $\log_a x = b$ ($a > 0, a \neq 1$)}
Phương trình luôn có \textbf{nghiệm duy nhất} với mọi số thực $b$:
$$\log_a x = b \iff x = a^b$$
\end{pedabox}

\begin{pedabox}{Phương pháp đưa về cùng cơ số}
Với $a > 0, a \neq 1$:
$$\log_a f(x) = \log_a g(x) \iff \begin{cases} f(x) > 0 \quad (\text{hoặc } g(x) > 0) \\ f(x) = g(x) \end{cases}$$
\end{pedabox}

\begin{scaffoldbox}{Biện pháp sư phạm: Bắt buộc đặt điều kiện xác định (Scaffolding)}
\textbf{Lưu ý cốt tử cho học sinh trung bình - yếu:}
\begin{itemize}
    \item \textbf{Quy tắc sống còn:} Khi giải phương trình lôgarit, bước đầu tiên LUÔN LUÔN là \textbf{đặt điều kiện xác định} cho biểu thức dưới dấu lôgarit phải dương: $f(x) > 0$ và $g(x) > 0$.
    \item Sau khi tìm được nghiệm, BẮT BUỘC phải đối chiếu lại với điều kiện để \textbf{loại bỏ nghiệm ngoại lai}!
\end{itemize}
\end{scaffoldbox}

\begin{aibox}{Tích hợp Năng lực AI (Theo Quyết định 2422/QĐ-BGDĐT - Mã 11.A1.1)}
\textbf{Phản biện và kiểm chứng lời giải AI:}
\begin{itemize}[leftmargin=0.6cm]
    \item \textbf{Tình huống thực tế:} Khi hỏi một ứng dụng AI giải phương trình $\log_3(x + 1) = \log_3(x^2 - 1)$, AI thực hiện:
    $$x + 1 = x^2 - 1 \iff x^2 - x - 2 = 0 \iff \left[\begin{array}{l} x = -1 \\ x = 2 \end{array}\right.$$
    AI kết luận phương trình có 2 nghiệm là $x = -1$ và $x = 2$.
    \item \textbf{Học sinh đánh giá:} Lời giải của AI bị \textbf{SAI NGHIÊM TRỌNG}!
    \item \textbf{Giải thích:} Với $x = -1$, ta có $x + 1 = 0 \implies \log_3 0$ không xác định.
    Điều kiện xác định là: $\begin{cases} x + 1 > 0 \\ x^2 - 1 > 0 \end{cases} \iff x > 1$. Do đó nghiệm $x = -1$ bị \textbf{LOẠI}, phương trình chỉ có duy nhất $1$ nghiệm là $x = 2$.
    \item \textbf{Bài học năng lực AI:} Học sinh không được tin tưởng mù quáng vào kết quả của AI mà phải luôn kiểm chứng tính chặt chẽ về điều kiện toán học.
\end{itemize}
\end{aibox}

\noindent \textbf{c) Sản phẩm:}
Học sinh hoàn thành lời giải bài toán mẫu:
\begin{itemize}
    \item \textbf{Ví dụ 1:} Giải phương trình $\log_2 (x + 1) = 3$.
    \begin{itemize}
        \item Điều kiện: $x + 1 > 0 \iff x > -1$.
        \item Phương trình tương đương: $x + 1 = 2^3 = 8 \iff x = 7$.
        \item Đối chiếu điều kiện: $x = 7 > -1$ (thỏa mãn). Vậy tập nghiệm là $S = \{7\}$.
    \end{itemize}
    \item \textbf{Ví dụ 2:} Giải phương trình $\log_3 (x + 1) = \log_3 (x^2 - 1)$.
    \begin{itemize}
        \item Điều kiện: $\begin{cases} x + 1 > 0 \\ x^2 - 1 > 0 \end{cases} \iff x > 1$.
        \item Phương trình tương đương: $x + 1 = x^2 - 1 \iff x^2 - x - 2 = 0 \iff \left[\begin{array}{l} x = -1 \\ x = 2 \end{array}\right.$.
        \item Đối chiếu điều kiện: $x = -1$ (loại); $x = 2$ (thỏa mãn). Vậy tập nghiệm là $S = \{2\}$.
    \end{itemize}
\end{itemize}

\noindent \textbf{d) Tổ chức thực hiện:}
\begin{itemize}[leftmargin=0.6cm]
    \item \textbf{Giao nhiệm vụ:} Giáo viên trình bày định nghĩa phương trình lôgarit; trình chiếu tình huống phản biện lỗi sai của AI. Yêu cầu học sinh giải Ví dụ 1 và 2.
    \item \textbf{Thực hiện nhiệm vụ:} Học sinh làm bài vào vở, chú ý ghi rõ bước điều kiện và bước đối chiếu điều kiện.
    \item \textbf{Báo cáo, thảo luận:} 2 học sinh lên bảng trình bày; cả lớp phân tích tại sao loại $x = -1$.
    \item \textbf{Kết luận, nhận định:} Giáo viên nhấn mạnh: Thiếu bước điều kiện và đối chiếu nghiệm sẽ bị trừ điểm trực tiếp trong các bài thi tự luận.
\end{itemize}

\subsubsection*{3. Hoạt động 3: Luyện tập (10 phút)}

\noindent \textbf{a) Mục tiêu:}
Rèn luyện kỹ năng giải thành thạo các phương trình mũ và lôgarit cơ bản; sử dụng MTCT để thử lại nghiệm.

\noindent \textbf{b) Nội dung:}
Thực hiện các bài tập trong Phiếu học tập số 1:
\begin{enumerate}[leftmargin=0.6cm]
    \item \textbf{Bài 1:} Giải các phương trình mũ sau:
    $$a) \ 5^{2x - 3} = 125; \qquad b) \ \left(\dfrac{1}{3}\right)^{x^2 - 2x} = 27$$
    \item \textbf{Bài 2:} Giải các phương trình lôgarit sau:
    $$a) \ \log_5 (2x - 1) = 2; \qquad b) \ \log_2 (x - 2) + \log_2 (x - 3) = 1$$
\end{enumerate}

\noindent \textbf{c) Sản phẩm:}
Lời giải chi tiết:
\begin{itemize}
    \item \textbf{Bài 1:}
    \begin{itemize}
        \item $a)$ $5^{2x - 3} = 5^3 \iff 2x - 3 = 3 \iff 2x = 6 \iff x = 3$. Tập nghiệm $S = \{3\}$.
        \item $b)$ $\left(3^{-1}\right)^{x^2 - 2x} = 3^3 \iff 3^{-x^2 + 2x} = 3^3 \iff -x^2 + 2x = 3 \iff x^2 - 2x + 3 = 0$ (vô nghiệm vì $\Delta' = 1 - 3 = -2 < 0$). Tập nghiệm $S = \emptyset$.
    \end{itemize}
    \item \textbf{Bài 2:}
    \begin{itemize}
        \item $a)$ Điều kiện: $2x - 1 > 0 \iff x > \dfrac{1}{2}$.
        Phương trình $\iff 2x - 1 = 5^2 = 25 \iff 2x = 26 \iff x = 13$ (thỏa mãn). Tập nghiệm $S = \{13\}$.
        \item $b)$ Điều kiện: $\begin{cases} x - 2 > 0 \\ x - 3 > 0 \end{cases} \iff x > 3$.
        Phương trình $\iff \log_2 [(x - 2)(x - 3)] = 1 \iff (x - 2)(x - 3) = 2^1 = 2$.
        $\iff x^2 - 5x + 6 = 2 \iff x^2 - 5x + 4 = 0 \iff \left[\begin{array}{l} x = 1 \\ x = 4 \end{array}\right.$.
        Đối chiếu điều kiện $x > 3$: Loại $x = 1$, nhận $x = 4$. Tập nghiệm $S = \{4\}$.
    \end{itemize}
\end{itemize}

\noindent \textbf{d) Tổ chức thực hiện:}
\begin{itemize}[leftmargin=0.6cm]
    \item \textbf{Giao nhiệm vụ:} Học sinh làm việc cá nhân 6 phút.
    \item \textbf{Thực hiện nhiệm vụ:} Giáo viên đi quanh lớp quan sát, nhắc nhở học sinh ở Bài 2b phải gom tổng hai lôgarit thành lôgarit của một tích trước khi giải.
    \item \textbf{Báo cáo, thảo luận:} 2 học sinh lên bảng giải; lớp nhận xét.
    \item \textbf{Kết luận, nhận định:} Giáo viên chốt lại phương pháp giải phương trình lôgarit chứa tổng các lôgarit.
\end{itemize}

\subsubsection*{4. Hoạt động 4: Vận dụng (5 phút)}

\noindent \textbf{a) Mục tiêu:}
Vận dụng phương trình mũ giải bài toán tài chính tính thời gian đạt mục tiêu tiết kiệm.

\noindent \textbf{b) Nội dung:}
\begin{pedabox}{Bài toán Tài chính: Thời gian gấp đôi số vốn}
Bác An gửi tiết kiệm $200$ triệu đồng vào ngân hàng với lãi suất $7\%$/năm theo thể thức lãi kép hàng năm. Công thức số tiền cả gốc lẫn lãi nhận được sau $n$ năm là:
$$A_n = 200 \cdot (1 + 0{,}07)^n = 200 \cdot (1{,}07)^n\text{ (triệu đồng)}$$
Hỏi sau ít nhất bao nhiêu năm thì bác An nhận được số tiền ít nhất là $400$ triệu đồng (tức là số tiền ban đầu tăng gấp đôi)?
\end{pedabox}

\noindent \textbf{c) Sản phẩm:}
Thiết lập phương trình:
$$200 \cdot (1{,}07)^n = 400 \iff (1{,}07)^n = 2 \iff n = \log_{1{,}07} 2$$
Sử dụng MTCT bấm: $n = \log_{1{,}07} 2 \approx 10{,}24\text{ năm}$.
Vì số kỳ hạn gửi tiền là từng năm (nguyên dương), nên sau ít nhất $11$ năm thì bác An mới nhận được số tiền vượt quá $400$ triệu đồng (cụ thể sau 11 năm nhận được $\approx 420{,}97$ triệu đồng).

\noindent \textbf{d) Tổ chức thực hiện:}
\begin{itemize}[leftmargin=0.6cm]
    \item \textbf{Giao nhiệm vụ:} Giáo viên giao bài toán, gợi ý học sinh lấy lôgarit cơ số $1{,}07$ của $2$.
    \item \textbf{Thực hiện nhiệm vụ:} Học sinh tính toán trên MTCT và biện luận số năm nguyên.
    \item \textbf{Báo cáo, nhận định:} Giáo viên biểu dương học sinh và giao nhiệm vụ chuẩn bị cho Tiết 2.
\end{itemize}

\vspace{10pt}

\subsection*{\faClockO \ TIẾT 2 (TIẾT 61 THEO PPCT): BẤT PHƯƠNG TRÌNH MŨ, BẤT PHƯƠNG TRÌNH LÔGARIT VÀ ỨNG DỤNG THỰC TIỄN}

\subsubsection*{1. Hoạt động 1: Mở đầu (Khởi động) (5 phút)}

\noindent \textbf{a) Mục tiêu:}
Tạo mâu thuẫn nhận thức về tính đổi chiều của bất phương trình khi cơ số nhỏ hơn 1 thông qua so sánh lũy thừa.

\noindent \textbf{b) Nội dung:}
Giáo viên đặt vấn đề:
\begin{pedabox}{Tình huống khởi động: Thử thách so sánh dấu}
\begin{enumerate}[leftmargin=0.6cm]
    \item Cho $2^x > 8$. Vì $2^3 = 8$ và cơ số $2 > 1$, ta suy ra điều gì về $x$ và số $3$?
    \item Cho $\left(\dfrac{1}{2}\right)^x > \dfrac{1}{8}$. Vì $\left(\dfrac{1}{2}\right)^3 = \dfrac{1}{8}$, ta suy ra $x > 3$ hay $x < 3$?
\end{enumerate}
\end{pedabox}

\noindent \textbf{c) Sản phẩm:}
Câu trả lời của học sinh:
\begin{itemize}
    \item Với $2^x > 8 = 2^3$: Do hàm số $y = 2^x$ đồng biến ($a = 2 > 1$) nên $x > 3$ (giữ nguyên chiều).
    \item Với $\left(\dfrac{1}{2}\right)^x > \dfrac{1}{8} = \left(\dfrac{1}{2}\right)^3$: Do hàm số $y = \left(\dfrac{1}{2}\right)^x$ nghịch biến ($a = \dfrac{1}{2} < 1$) nên $x < 3$ (\textbf{bắt buộc đổi chiều}!).
    \textit{Thử lại:} Với $x = 2 < 3$, ta có $\left(\dfrac{1}{2}\right)^2 = \dfrac{1}{4} = \dfrac{2}{8} > \dfrac{1}{8}$ (đúng!).
\end{itemize}

\noindent \textbf{d) Tổ chức thực hiện:}
\begin{itemize}[leftmargin=0.6cm]
    \item \textbf{Giao nhiệm vụ:} Giáo viên nêu câu hỏi, yêu cầu học sinh thảo luận cặp đôi trong 2 phút.
    \item \textbf{Thực hiện nhiệm vụ:} Học sinh đối chiếu tính chất đồng biến/nghịch biến của hàm số mũ ở Bài 20 để rút ra nhận xét.
    \item \textbf{Báo cáo, nhận định:} Giáo viên dẫn dắt: Khi giải bất phương trình mũ và lôgarit, việc quan sát cơ số lớn hơn 1 hay nhỏ hơn 1 là yếu tố quyết định chiều của bất phương trình.
\end{itemize}

\subsubsection*{2. Hoạt động 2: Hình thành kiến thức (25 phút)}

\paragraph{2.1. Đơn vị kiến thức 1: Bất phương trình mũ cơ bản (12 phút)}

\noindent \textbf{a) Mục tiêu:}
Nắm vững quy tắc giải bất phương trình mũ cơ bản; giải thích được tại sao phải đổi chiều khi cơ số thuộc khoảng $(0; 1)$.

\noindent \textbf{b) Nội dung:}
\begin{pedabox}{Bất phương trình mũ cơ bản}
Cho số thực dương $a \neq 1$.
\begin{itemize}[leftmargin=0.6cm]
    \item \textbf{Trường hợp cơ số $a > 1$ (Giữ nguyên chiều):}
    \begin{itemize}
        \item $a^{f(x)} > a^{g(x)} \iff f(x) > g(x)$.
        \item $a^x > b$: Nếu $b \le 0$, BPT nghiệm đúng với mọi $x \in \mathbb{R}$; nếu $b > 0$, BPT $\iff x > \log_a b$.
    \end{itemize}
    \item \textbf{Trường hợp cơ số $0 < a < 1$ (Bắt buộc ĐỔI CHIỀU):}
    \begin{itemize}
        \item $a^{f(x)} > a^{g(x)} \iff f(x) < g(x)$.
        \item $a^x > b$: Nếu $b \le 0$, BPT nghiệm đúng với mọi $x \in \mathbb{R}$; nếu $b > 0$, BPT $\iff x < \log_a b$.
    \end{itemize}
\end{itemize}
\end{pedabox}

\begin{scaffoldbox}{Giàn giáo sư phạm: Khẩu quyết giải bất phương trình mũ}
\textbf{Khẩu quyết vàng cho học sinh trung bình - yếu:}
\begin{center}
\textit{``Cơ số LỚN HƠN 1: GIỮ NGUYÊN CHIỀU; Cơ số NHỎ HƠN 1: ĐỔI CHIỀU BẤT ĐẲNG THỨC''}
\end{center}
\end{scaffoldbox}

\noindent \textbf{c) Sản phẩm:}
Học sinh hoàn thành các ví dụ áp dụng:
\begin{itemize}
    \item \textbf{Ví dụ 1:} Giải BPT $3^{2x - 1} \ge 9$.
    Vì cơ số $a = 3 > 1$, giữ nguyên chiều:
    $$3^{2x - 1} \ge 3^2 \iff 2x - 1 \ge 2 \iff 2x \ge 3 \iff x \ge \dfrac{3}{2}.$$
    Tập nghiệm là $S = \left[\dfrac{3}{2}; +\infty\right)$.
    \item \textbf{Ví dụ 2:} Giải BPT $\left(\dfrac{1}{2}\right)^{x+1} > 4$.
    Ta có $4 = 2^2 = \left(\dfrac{1}{2}\right)^{-2}$.
    Vì cơ số $a = \dfrac{1}{2} \in (0; 1)$, đổi chiều:
    $$\left(\dfrac{1}{2}\right)^{x+1} > \left(\dfrac{1}{2}\right)^{-2} \iff x + 1 < -2 \iff x < -3.$$
    Tập nghiệm là $S = (-\infty; -3)$.
\end{itemize}

\noindent \textbf{d) Tổ chức thực hiện:}
\begin{itemize}[leftmargin=0.6cm]
    \item \textbf{Giao nhiệm vụ:} Giáo viên trình bày quy tắc, hướng dẫn Ví dụ 1 và yêu cầu học sinh làm Ví dụ 2.
    \item \textbf{Thực hiện nhiệm vụ:} Học sinh làm bài vào vở, chú ý bước đổi chiều ở Ví dụ 2.
    \item \textbf{Báo cáo, thảo luận:} Gọi 1 học sinh lên bảng giải Ví dụ 2; cả lớp nhận xét bước đổi chiều.
    \item \textbf{Kết luận, nhận định:} Giáo viên nhấn mạnh: Cứ thấy cơ số là phân số nhỏ hơn 1 hoặc số thập phân như $0{,}5; 0{,}2$ thì phải lập tức đánh dấu đổi chiều.
\end{itemize}

\paragraph{2.2. Đơn vị kiến thức 2: Bất phương trình lôgarit cơ bản (13 phút)}

\noindent \textbf{a) Mục tiêu:}
Nắm vững quy tắc giải bất phương trình lôgarit cơ bản; luôn ghi nhớ kết hợp điều kiện xác định và quy tắc đổi chiều khi cơ số nhỏ hơn 1.

\noindent \textbf{b) Nội dung:}
\begin{pedabox}{Bất phương trình lôgarit cơ bản}
Cho số thực dương $a \neq 1$.
\begin{itemize}[leftmargin=0.6cm]
    \item \textbf{Trường hợp cơ số $a > 1$ (Giữ nguyên chiều):}
    $$\log_a f(x) > \log_a g(x) \iff \begin{cases} g(x) > 0 \\ f(x) > g(x) \end{cases}$$
    $$\log_a x > b \iff x > a^b$$
    $$\log_a x < b \iff 0 < x < a^b \quad (\text{Bắt buộc kẹp thêm điều kiện } x > 0!)$$
    \item \textbf{Trường hợp cơ số $0 < a < 1$ (Bắt buộc ĐỔI CHIỀU):}
    $$\log_a f(x) > \log_a g(x) \iff \begin{cases} f(x) > 0 \\ f(x) < g(x) \end{cases}$$
    $$\log_a x > b \iff 0 < x < a^b \quad (\text{Đổi chiều và kẹp điều kiện } x > 0!)$$
    $$\log_a x < b \iff x > a^b$$
\end{itemize}
\end{pedabox}

\begin{scaffoldbox}{Sai lầm kinh điển cần triệt tiêu trong BPT Lôgarit}
\textbf{Cảnh báo nghiêm khắc:}
\begin{itemize}
    \item Rất nhiều học sinh khi giải: $\log_2 x < 3 \iff x < 2^3 = 8$ và kết luận tập nghiệm là $(-\infty; 8) \implies$ \textbf{SAI HOÀN TOÀN}!
    \item Vì biểu thức trong dấu lôgarit phải dương, nên nghiệm đúng phải là: $0 < x < 8$, tập nghiệm là $(0; 8)$.
    \item Khi giải: $\log_{0{,}5} x > 2 \iff x < (0{,}5)^2 = 0{,}25$. Nhưng quên kẹp điều kiện $x > 0$ thì sẽ lấy cả các số âm $\implies$ \textbf{SAI}! Nghiệm đúng là: $0 < x < 0{,}25$.
\end{itemize}
\end{scaffoldbox}

\noindent \textbf{c) Sản phẩm:}
Học sinh hoàn thành các ví dụ:
\begin{itemize}
    \item \textbf{Ví dụ 1:} Giải BPT $\log_2 (2x - 4) \le 3$.
    \begin{itemize}
        \item Điều kiện: $2x - 4 > 0 \iff x > 2$.
        \item Vì cơ số $a = 2 > 1$, giữ nguyên chiều:
        $$2x - 4 \le 2^3 = 8 \iff 2x \le 12 \iff x \le 6.$$
        \item Kết hợp điều kiện: $2 < x \le 6$. Tập nghiệm là $S = (2; 6]$.
    \end{itemize}
    \item \textbf{Ví dụ 2:} Giải BPT $\log_{\frac{1}{3}} (x - 1) \ge -2$.
    \begin{itemize}
        \item Điều kiện: $x - 1 > 0 \iff x > 1$.
        \item Vì cơ số $a = \dfrac{1}{3} \in (0; 1)$, đổi chiều:
        $$x - 1 \le \left(\dfrac{1}{3}\right)^{-2} = 3^2 = 9 \iff x \le 10.$$
        \item Kết hợp điều kiện: $1 < x \le 10$. Tập nghiệm là $S = (1; 10]$.
    \end{itemize}
\end{itemize}

\noindent \textbf{d) Tổ chức thực hiện:}
\begin{itemize}[leftmargin=0.6cm]
    \item \textbf{Giao nhiệm vụ:} Giáo viên trình bày quy tắc, phân tích bẫy kẹp điều kiện $x > 0$. Yêu cầu học sinh làm Ví dụ 1 và 2.
    \item \textbf{Thực hiện nhiệm vụ:} Học sinh làm bài, đối chiếu điều kiện để tìm tập nghiệm.
    \item \textbf{Báo cáo, thảo luận:} 2 học sinh lên bảng giải; lớp nhận xét bước kết hợp nghiệm trên trục số.
    \item \textbf{Kết luận, nhận định:} Giáo viên nhấn mạnh: Luôn viết điều kiện thành một bất đẳng thức kép để không bao giờ bị sót điều kiện dương.
\end{itemize}

\subsubsection*{3. Hoạt động 3: Luyện tập (10 phút)}

\noindent \textbf{a) Mục tiêu:}
Rèn luyện kỹ năng giải bất phương trình mũ và lôgarit tổng hợp, giao các tập nghiệm chính xác.

\noindent \textbf{b) Nội dung:}
Thực hiện các bài tập trong Phiếu học tập số 2:
\begin{enumerate}[leftmargin=0.6cm]
    \item \textbf{Bài 1:} Giải các bất phương trình sau:
    $$a) \ 2^{x^2 - 3x} < 16; \qquad b) \ \left(\dfrac{1}{5}\right)^{2x - 1} \ge \left(\dfrac{1}{5}\right)^{x + 3}$$
    \item \textbf{Bài 2:} Giải bất phương trình:
    $$\log_3 (x + 2) + \log_3 (x - 4) \le 3$$
\end{enumerate}

\noindent \textbf{c) Sản phẩm:}
Lời giải chi tiết:
\begin{itemize}
    \item \textbf{Bài 1:}
    \begin{itemize}
        \item $a)$ $2^{x^2 - 3x} < 2^4 \iff x^2 - 3x < 4 \iff x^2 - 3x - 4 < 0 \iff -1 < x < 4$. Tập nghiệm $S = (-1; 4)$.
        \item $b)$ Vì cơ số $\dfrac{1}{5} < 1$, đổi chiều: $2x - 1 \le x + 3 \iff x \le 4$. Tập nghiệm $S = (-\infty; 4]$.
    \end{itemize}
    \item \textbf{Bài 2:}
    \begin{itemize}
        \item Điều kiện: $\begin{cases} x + 2 > 0 \\ x - 4 > 0 \end{cases} \iff x > 4$.
        \item BPT $\iff \log_3 [(x + 2)(x - 4)] \le 3 \iff (x + 2)(x - 4) \le 3^3 = 27$ (vì cơ số $3 > 1$).
        $\iff x^2 - 2x - 8 \le 27 \iff x^2 - 2x - 35 \le 0 \iff -5 \le x \le 7$.
        \item Kết hợp điều kiện $x > 4$: Ta được $4 < x \le 7$.
        Tập nghiệm là $S = (4; 7]$.
    \end{itemize}
\end{itemize}

\noindent \textbf{d) Tổ chức thực hiện:}
\begin{itemize}[leftmargin=0.6cm]
    \item \textbf{Giao nhiệm vụ:} Học sinh làm bài độc lập trong 6 phút.
    \item \textbf{Thực hiện nhiệm vụ:} Giáo viên quan sát, giúp đỡ học sinh yếu bước giải bất phương trình bậc hai tam thức bậc hai.
    \item \textbf{Báo cáo, thảo luận:} 2 học sinh lên bảng giải; cả lớp nhận xét.
    \item \textbf{Kết luận, nhận định:} Giáo viên chuẩn hóa các bước kết hợp nghiệm trên trục số.
\end{itemize}

\subsubsection*{4. Hoạt động 4: Vận dụng (5 phút)}

\noindent \textbf{a) Mục tiêu:}
Vận dụng bất phương trình lôgarit giải quyết bài toán đo lường độ chấn động động đất theo thang đo Richter.

\noindent \textbf{b) Nội dung:}
\begin{stembox}{Bài toán Khoa học Trái đất: Độ chấn động Richter của các trận động đất}
Độ lớn $M$ của một trận động đất theo thang đo Richter được tính bằng công thức:
$$M = \log A - \log A_0 = \log\left(\dfrac{A}{A_0}\right)$$
(trong đó $A$ là biên độ tối đa đo được bởi máy địa chấn và $A_0$ là biên độ chuẩn của một trận động đất rất nhẹ ở cùng khoảng cách).
Một khu vực xây dựng công trình yêu cầu nền móng phải chịu được các trận động đất có độ chấn động không vượt quá $6{,}5$ độ Richter ($M \le 6{,}5$).
Hỏi biên độ chấn động tối đa $A$ của khu vực này được phép gấp tối đa bao nhiêu lần biên độ chuẩn $A_0$?
\end{stembox}

\noindent \textbf{c) Sản phẩm:}
Thiết lập bất phương trình:
$$M \le 6{,}5 \iff \log\left(\dfrac{A}{A_0}\right) \le 6{,}5$$
Vì đây là lôgarit thập phân cơ số $10 > 1$, ta giữ nguyên chiều:
$$\dfrac{A}{A_0} \le 10^{6{,}5} \approx 3\,162\,277{,}66$$
Vậy biên độ chấn động tối đa $A$ được phép gấp khoảng $3{,}16$ triệu lần biên độ chuẩn $A_0$.

\noindent \textbf{d) Tổ chức thực hiện:}
\begin{itemize}[leftmargin=0.6cm]
    \item \textbf{Giao nhiệm vụ:} Giáo viên giao bài toán thực tế, yêu cầu học sinh chuyển đổi về bất phương trình lôgarit cơ số 10.
    \item \textbf{Thực hiện nhiệm vụ:} Học sinh tính lũy thừa $10^{6{,}5}$ trên MTCT.
    \item \textbf{Báo cáo, nhận định:} Học sinh phát biểu kết quả; giáo viên tổng kết toàn bài và hướng dẫn chuẩn bị cho bài ôn tập cuối chương VI.
\end{itemize}

\vspace{10pt}

\section*{IV. PHỤ LỤC HỌC LIỆU BỔ TRỢ (BẮT BUỘC)}

\subsection*{1. Phiếu học tập số 1 (Tiết 1): Phương trình mũ và Phương trình lôgarit cơ bản}

\begin{pedabox}{PHIẾU HỌC TẬP SỐ 1: GIẢI PHƯƠNG TRÌNH MŨ VÀ LÔGARIT}
\textbf{Họ và tên học sinh:} \dotfill \ \textbf{Lớp:} \dotfill \ \textbf{Nhóm:} \dotfill

\noindent \textbf{CẤP ĐỘ 1: NHẬN BIẾT (Khởi động nhanh)}
\begin{enumerate}[leftmargin=0.6cm]
    \item Giải các phương trình mũ và lôgarit cơ bản sau:
    $$a) \ 2^x = 16 \iff x = \dots\dots; \qquad b) \ 3^{x-1} = \dfrac{1}{9} \iff x = \dots\dots; \qquad c) \ \log_3 x = 4 \iff x = \dots\dots$$
    \item Tìm điều kiện xác định của các phương trình sau (chưa cần giải):
    $$a) \ \log_2 (3x - 6) = 1; \qquad b) \ \log_5 (x^2 - 9) = 2; \qquad c) \ \log_3 (x - 1) = \log_3 (2x - 5)$$
\end{enumerate}

\noindent \textbf{CẤP ĐỘ 2: THÔNG HIỂU (Đưa về cùng cơ số \& Bấm MTCT kiểm tra)}
\begin{enumerate}[leftmargin=0.6cm]
    \item Giải phương trình: $4^{x+1} = 8^{2x - 1}$.
    \item Giải phương trình: $\log_2 (x + 3) + \log_2 (x - 1) = 2$.
    \item Sử dụng tính năng \texttt{SOLVE} trên MTCT Casio fx-580VN X để tìm nghiệm phương trình $3^{x} = 10$.
\end{enumerate}

\noindent \textbf{CẤP ĐỘ 3: VẬN DỤNG (Phản biện lời giải AI)}
\begin{enumerate}[leftmargin=0.6cm]
    \item Một AI giải phương trình $\log_2(x^2 - 4) = \log_2(3x)$ ra hai nghiệm $x = -1$ và $x = 4$. Hãy chỉ ra lỗi sai của AI và viết lại lời giải chuẩn xác.
\end{enumerate}
\end{pedabox}

\subsection*{2. Phiếu học tập số 2 (Tiết 2): Bất phương trình mũ, lôgarit và Ứng dụng}

\begin{pedabox}{PHIẾU HỌC TẬP SỐ 2: BẤT PHƯƠNG TRÌNH MŨ VÀ LÔGARIT}
\textbf{Họ và tên học sinh:} \dotfill \ \textbf{Lớp:} \dotfill \ \textbf{Nhóm:} \dotfill

\noindent \textbf{CẤP ĐỘ 1: NHẬN BIẾT \& THÔNG HIỂU (Nhận diện cơ số \& Đổi chiều)}
\begin{enumerate}[leftmargin=0.6cm]
    \item Điền dấu thích hợp ($>, <, \ge, \le$) vào chỗ trống:
    $$a) \ 2^x > 4 \iff x \ \dots\dots \ 2; \qquad b) \ (0{,}3)^x > 0{,}09 \iff x \ \dots\dots \ 2; \qquad c) \ \log_{\frac{1}{2}} x \ge 3 \iff x \ \dots\dots \ \dfrac{1}{8}$$
    \item Giải các bất phương trình sau:
    $$a) \ 5^{x - 2} \le 125; \qquad b) \ \log_3 (2x - 1) < 2; \qquad c) \ \log_{0{,}2} (x + 1) \ge -1$$
\end{enumerate}

\noindent \textbf{CẤP ĐỘ 2: VẬN DỤNG (Bất phương trình chứa biến phức tạp)}
\begin{enumerate}[leftmargin=0.6cm]
    \item Giải bất phương trình: $\left(\dfrac{1}{2}\right)^{x^2 - 5x + 6} \ge 1$.
    \item Giải bất phương trình: $\log_2 (x - 1) - \log_2 (x + 2) \le 1$.
\end{enumerate}

\noindent \textbf{CẤP ĐỘ 3: VẬN DỤNG CAO (Mô hình hóa thực tiễn)}
\begin{enumerate}[leftmargin=0.6cm]
    \item Độ pH của dung dịch được tính bởi $\text{pH} = -\log[\text{H}^+]$. Một hồ bơi an toàn khi độ pH thỏa mãn $7{,}2 \le \text{pH} \le 7{,}6$. Hãy giải bất phương trình kép để tìm khoảng nồng độ mol ion $[\text{H}^+]$ an toàn của hồ bơi.
\end{enumerate}
\end{pedabox}

\subsection*{3. Bảng Rubric đánh giá năng lực học sinh theo 4 mức độ}

\begin{center}
\begin{tabularx}{\textwidth}{|p{2.8cm}|X|X|X|X|}
    \hline
    \textbf{Tiêu chí} & \textbf{Mức 1: Chưa đạt} & \textbf{Mức 2: Đạt} & \textbf{Mức 3: Khá} & \textbf{Mức 4: Tốt} \\ \hline
    \textbf{1. Kỹ năng giải phương trình mũ, lôgarit} & Thường xuyên quên điều kiện xác định; lấy cả nghiệm âm của hàm mũ. & Nhớ đặt điều kiện cho biểu thức lôgarit; giải được phương trình cơ bản. & Giải thành thạo phương trình đưa về cùng cơ số; biết gom lôgarit cùng cơ số. & Lập luận sắc bén, biến đổi linh hoạt; giải nhanh các phương trình quy về bậc hai. \\ \hline
    \textbf{2. Kỹ năng giải bất phương trình \& Đổi chiều} & Quên đổi chiều khi cơ số nhỏ hơn 1; quên kẹp điều kiện dương cho BPT lôgarit. & Biết đổi chiều khi cơ số $a < 1$; giải được BPT mũ, lôgarit cơ bản. & Luôn kẹp đúng điều kiện dương; kết hợp nghiệm trên trục số chính xác. & Xử lý điêu luyện bất phương trình chứa tham số và hệ bất phương trình mũ, lôgarit. \\ \hline
    \textbf{3. Năng lực số \& Dò nghiệm MTCT (Mã 5.2NC1b)} & Chưa biết nhập phương trình và bấm \texttt{SOLVE}; không biết đọc kết quả. & Bấm được lệnh \texttt{SOLVE} theo hướng dẫn để tìm nghiệm phương trình đơn giản. & Thao tác bấm máy nhanh; biết đổi giá trị khởi tạo \texttt{X=?} để dò nghiệm khác. & Làm chủ tính năng \texttt{SOLVE} và \texttt{TABLE} để thử lại nghiệm và dấu của BPT. \\ \hline
    \textbf{4. Đánh giá lời giải AI \& Vận dụng (Mã 11.A1.1)} & Tin tưởng tuyệt đối kết quả của AI mà không kiểm tra lại điều kiện nghiệm. & Nhận ra lời giải AI thiếu điều kiện khi được giáo viên gợi ý. & Tự phát hiện và phân tích chính xác lỗi nghiệm ngoại lai của AI; giải thích được. & Xây dựng prompt AI chuẩn mực; vận dụng toán mũ, lôgarit giải bài toán thực tiễn xuất sắc. \\ \hline
\end{tabularx}
\end{center}

\end{document}
